\thispagestyle{plain}
\noindent
\begin{tikzpicture}
\node[fill=fnDark,text=white,rounded corners=2pt,inner sep=6pt,
      minimum width=\textwidth,anchor=west]
  {\large\bfseries Pinout summary --- pin-by-pin connection};
\end{tikzpicture}

\vspace{6pt}
\noindent
{\footnotesize
Quick-reference table: the \textbf{57 real signals} of the top-level
\code{spi\_neuron\_top}, each with its own individual \code{CABGA381} ball
(\textbf{not} a bus range) --- real data from Project~Trellis's device
database (\code{iodb.json}), \textbf{verified by a complete
\code{nextpnr-ecp5} place\&route run at 0 errors} (not a planned pinout).
Full description, per-bank placement rationale and the pin-by-pin
connection to the ISSI PSRAM: ch.~\ref{ch:hw}.
}

\vspace{4pt}
\noindent
\renewcommand{\arraystretch}{1.08}
\begin{tabularx}{\textwidth}{L{2.7cm} C{1.0cm} C{1.0cm} C{0.9cm} Y}
\toprule
\rowh \thd{Signal} & \thd{Ball} & \thd{Bank} & \thd{Dir} & \thd{Corresponding pin / function} \\
\midrule
\multicolumn{5}{l}{\textit{\color{fnDark}Clock and reset}}\\
\code{clk}  & H5 & 7 & IN  & System clock, pad \code{GR\_PCLK7\_0} (dedicated global clock). \\
\rowa \code{rst} & B4 & 7 & IN & Global synchronous reset, active high. \\
\multicolumn{5}{l}{\textit{\color{fnDark}Application SPI (host $\leftrightarrow$ FPGA, Mode~0)}}\\
\code{sclk} & B5 & 7 & IN  & SPI clock (CPOL=0, CPHA=0). \\
\rowa \code{mosi} & C5 & 7 & IN & Master-Out Slave-In. \\
\code{miso} & A3 & 7 & OUT & Master-In Slave-Out. \\
\rowa \code{cs\_n} & B3 & 7 & IN & Chip-select, active low. \\
\multicolumn{5}{l}{\textit{\color{fnDark}Host attention (active-low, level)}}\\
\code{data\_ready\_n} & C3 & 7 & OUT & Low while a result is waiting to be read. \\
\rowa \code{irq\_n} & C4 & 7 & OUT & Low while the graph engine's load-time guard has tripped. \\
\multicolumn{5}{l}{\textit{\color{fnDark}Flash subsystem --- SPI toward W25Q128JV (boot/persistence)}}\\
\code{flash\_sclk} & E3 & 7 & OUT & SPI clock toward the flash --- ordinary GPIO, independent (Phase F7, ch.~\ref{ch:hw}). \\
\rowa \code{flash\_mosi} & D3 & 7 & OUT & Master-Out Slave-In toward the onboard flash. \\
\code{flash\_miso} & D5 & 7 & IN & Master-In Slave-Out from the flash. \\
\rowa \code{flash\_cs\_n} & E4 & 7 & OUT & Flash chip-select, active low. \\
\multicolumn{5}{l}{\textit{\color{fnDark}PSRAM address bus \code{psram\_a[21:0]} --- 22 individual balls (bank 2)}}\\
\code{psram\_a[0]}  & E16 & 2 & OUT & PSRAM A0 \\
\rowa \code{psram\_a[1]}  & F16 & 2 & OUT & PSRAM A1 \\
\code{psram\_a[2]}  & D18 & 2 & OUT & PSRAM A2 \\
\rowa \code{psram\_a[3]}  & E17 & 2 & OUT & PSRAM A3 \\
\code{psram\_a[4]}  & E18 & 2 & OUT & PSRAM A4 \\
\rowa \code{psram\_a[5]}  & F18 & 2 & OUT & PSRAM A5 \\
\code{psram\_a[6]}  & F17 & 2 & OUT & PSRAM A6 \\
\rowa \code{psram\_a[7]}  & G16 & 2 & OUT & PSRAM A7 \\
\code{psram\_a[8]}  & G18 & 2 & OUT & PSRAM A8 \\
\rowa \code{psram\_a[9]}  & H16 & 2 & OUT & PSRAM A9 \\
\code{psram\_a[10]} & H17 & 2 & OUT & PSRAM A10 \\
\rowa \code{psram\_a[11]} & H18 & 2 & OUT & PSRAM A11 \\
\code{psram\_a[12]} & J16 & 2 & OUT & PSRAM A12 \\
\rowa \code{psram\_a[13]} & J17 & 2 & OUT & PSRAM A13 \\
\code{psram\_a[14]} & C20 & 2 & OUT & PSRAM A14 \\
\rowa \code{psram\_a[15]} & D19 & 2 & OUT & PSRAM A15 \\
\code{psram\_a[16]} & E19 & 2 & OUT & PSRAM A16 \\
\rowa \code{psram\_a[17]} & E20 & 2 & OUT & PSRAM A17 \\
\code{psram\_a[18]} & F19 & 2 & OUT & PSRAM A18 \\
\rowa \code{psram\_a[19]} & F20 & 2 & OUT & PSRAM A19 \\
\code{psram\_a[20]} & G20 & 2 & OUT & PSRAM A20 \\
\rowa \code{psram\_a[21]} & H20 & 2 & OUT & PSRAM A21 \\
\code{psram\_a[22]} & P18 & 3 & OUT & Always 0 (byte$\to$word shift): NC on the board. \\
\multicolumn{5}{l}{\textit{\color{fnDark}PSRAM data bus \code{psram\_dq[15:0]} --- 16 individual balls (banks 2 and 3)}}\\
\rowa \code{psram\_dq[0]}  & K18 & 2 & IO & PSRAM DQ0 \\
\code{psram\_dq[1]}  & C18 & 2 & IO & PSRAM DQ1 (dual-function ball, used as ordinary GPIO). \\
\rowa \code{psram\_dq[2]}  & D17 & 2 & IO & PSRAM DQ2 \\
\code{psram\_dq[3]}  & D20 & 2 & IO & PSRAM DQ3 \\
\rowa \code{psram\_dq[4]}  & G19 & 2 & IO & PSRAM DQ4 \\
\code{psram\_dq[5]}  & J18 & 2 & IO & PSRAM DQ5 \\
\rowa \code{psram\_dq[6]}  & J19 & 2 & IO & PSRAM DQ6 \\
\code{psram\_dq[7]}  & J20 & 2 & IO & PSRAM DQ7 \\
\rowa \code{psram\_dq[8]}  & K19 & 2 & IO & PSRAM DQ8 \\
\code{psram\_dq[9]}  & K20 & 2 & IO & PSRAM DQ9 \\
\rowa \code{psram\_dq[10]} & L17 & 3 & IO & PSRAM DQ10 \\
\code{psram\_dq[11]} & M18 & 3 & IO & PSRAM DQ11 \\
\rowa \code{psram\_dq[12]} & M17 & 3 & IO & PSRAM DQ12 \\
\code{psram\_dq[13]} & N16 & 3 & IO & PSRAM DQ13 \\
\rowa \code{psram\_dq[14]} & N18 & 3 & IO & PSRAM DQ14 \\
\code{psram\_dq[15]} & P17 & 3 & IO & PSRAM DQ15 \\
\multicolumn{5}{l}{\textit{\color{fnDark}PSRAM control}}\\
\rowa \code{psram\_ce\_n} & N17 & 3 & OUT & PSRAM CE\# --- chip enable, active low. \\
\code{psram\_oe\_n} & R16 & 3 & OUT & PSRAM OE\# --- output enable (read). \\
\rowa \code{psram\_we\_n} & R17 & 3 & OUT & PSRAM WE\# --- write enable. \\
\code{psram\_lb\_n} & T16 & 3 & OUT & PSRAM LB\# --- lower-byte enable (DQ[7:0]). \\
\rowa \code{psram\_ub\_n} & N19 & 3 & OUT & PSRAM UB\# --- upper-byte enable (DQ[15:8]). \\
\code{psram\_zz\_n} & N20 & 3 & OUT & PSRAM ZZ\# --- sleep/snooze (high during normal operation). \\
\bottomrule
\end{tabularx}
\renewcommand{\arraystretch}{1.25}

\vspace{4pt}
\noindent
{\footnotesize\color{fnGrey}
Standard I/O: LVCMOS33 on all 57 signals. Boot config-SPI and JTAG balls (fixed-function
dedicated pins, no RTL port) do not appear in this table --- see ch.~\ref{ch:hw}
§``Configuration and programming''. Source: \code{synth/ecp5/spi\_neuron\_top.lpf},
generated by \code{tools/pinout/gen\_lpf.py} against Project~Trellis's
\code{iodb.json}.\par}
